\section{21--22 September 2026: Five adopted earlier-floor estimates}

All five remaining land corrections in the reviewed 23-case batch are applied
following Jacob's approval. Their recorded ground boundaries are supported;
earlier building-floor measures use explicitly flagged estimates or a later
administrative proxy. Exactly five parent rows change from the start of this
review. All constituent records, housing units, filing dates, parent memberships,
and the 554 historical/310 post-policy weighting sample remain unchanged.

East 125th uses 50,836.28 square feet of estimated earlier floor on 42,540
square feet of ground. Each old parcel's frozen floor is multiplied by its
own geometric overlap share, assuming uniform floor density within that parcel.
St. James uses the first post-split housing parcel's administrative 3,650 square
feet on 17,775 ground. Housing plus retained church floor is 1,700 below the
frozen old total; this discrepancy remains documented. Onderdonk uses
$18,658\times9,310/13,218=13,141.62$ square feet, a proportional allocation
whose uncertainty includes the location of the earlier second story.

Jamaica uses 39,349.5 square feet of residential ground and 40,342.49 of earlier
floor. The 2026 recorded declaration supplies the boundary, and 2024 condominium
plans locate shop floors, rear mezzanines, the second story, and the roof stair.
For each old parcel, its above-cellar plan share inside the boundary allocates
its frozen 2023 administrative floor total. Two shops outside the drawings use
recorded building depths. Six old parcels contribute housing ground; old lots
89 and 94 are wholly outside. Their 23,175 ground and 17,220 floor remain in
the eight-parcel audit inventory, while production counts six. Approximate shapes, scaled mezzanine depths, and
similar spatial allocation of unmeasured cellar floor remain assumptions.

Kingsbrook uses 105,382 square feet of ground: 92,709 for Phase I plus 12,673
of parking. HCR's corrected 2018 schedule gives 214,280 gross square feet for
four pavilions, including basements. The June 2025 environmental assessment
records an 11,600-square-foot power plant with a cellar. A tracing of the July
2021 survey places 47.9165 percent of that footprint inside Phase I, giving
5,558.31 per floor. Assuming an equally extensive cellar and a half-filled
mezzanine bay adds 662.37, producing 226,059.00 gross square feet. The mezzanine
midpoint is an imputation; neither its actual coverage nor the cellar partition
is measured. Omitting cellars gives 185,255.69. HCR gross area is a documented
substitute whose exact comparability with PLUTO is unproved. The source table's
excluded-floor amount is consequently an accounting residual, not a measured
retained-campus total.

The combined corrections change the weighted historical exact-99 share from
1.354308 to 1.353533 percent and the weighted mean from 151.4872 to 151.4686
units. The final calibration moment error is $6.37\times10^{-10}$. A separate
20-row sensitivity varies each site's earlier floor from zero to its complete
frozen old total while holding the other parents and formula fixed. These are
stress tests, not confidence intervals or joint bounds, and do not establish
insensitivity of the future structural model. The six flagged parents include
the previously approved Flatbush estimate. All 23 land corrections are applied;
GO Broome's dated earlier-floor source conflict remains queued separately.

\textit{Sources and reproduction.} The parcel audit's five-floor decision note
links the original plans and instruments with physical page references. Root
\texttt{make} rebuilds the panels and exhibits; \texttt{make five-floor-review}
reproduces the allocation worksheets, survey illustration, and sensitivity.
The stack-parcel audit now distinguishes the automatic filing lookup from
adopted parcel allocations, checking equality only for automatic sites.
Official new downloads are checksum-pinned. Production reads the manual source
table and archived parcel data; no main task depends on the audit.
